% Análisis matricial de armaduras con Python
% Compilar desde esta carpeta con: latexmk libro.tex
\documentclass[
  fontsize=10.5pt,
  twoside=true,
  open=right,
  headings=standardclasses,
  numbers=noenddot,
  bibliography=totoc,
]{scrbook}

\usepackage{estilo/armaduras}
\addbibresource{bibliografia.bib}

\title{Análisis matricial de armaduras con Python}
\subtitle{Aprende el método de rigidez construyendo tu propia calculadora}
\author{Eddy Eder Sucapuca Cruz}
\date{Versión preliminar}

\begin{document}

\frontmatter
% portada.tex: portada y página legal del libro.
% La armadura es el puente Pratt del ejemplo C (capítulo 12), con cada barra en
% el color de su fuerza axial: tracción, compresión o sin fuerza.

\begin{titlepage}
\begin{tikzpicture}[remember picture, overlay]
  % Franja superior con el título
  \fill[principal] (current page.north west) rectangle
    ([yshift=-0.52\paperheight]current page.north east);
  \fill[acento] ([yshift=-0.52\paperheight]current page.north west) rectangle
    ([yshift=-0.535\paperheight]current page.north east);
  \node[anchor=north west, text width=0.78\paperwidth, align=left,
      text=white, font=\sffamily]
    at ([xshift=0.11\paperwidth, yshift=-0.16\paperheight]current page.north west)
    {{\fontsize{34}{40}\selectfont\bfseries\spaceskip=0.28em
      Análisis matricial\\[2pt]
      de armaduras\\[2pt] con Python\par}
     \vspace{0.9cm}
     {\Large Aprende el método de rigidez\\[2pt]
      construyendo tu propia calculadora\par}};

  % El puente del ejemplo C
  % 9 unidades de luz con escala 1.3: centrado, 5.85 cm a cada lado del eje
  \begin{scope}[shift={([xshift=-5.85cm, yshift=-0.79\paperheight]
      current page.north)}, scale=1.3, line cap=round]
    \foreach \k in {0,...,6}
      \coordinate (p\the\numexpr\k+1\relax) at (1.5*\k, 0);
    \foreach \k in {0,...,4}
      \coordinate (p\the\numexpr\k+8\relax) at (1.5*\k+1.5, 1.5);
    \foreach \i/\j/\tono in {
        1/2/traccion, 2/3/traccion, 3/4/traccion, 4/5/traccion,
        5/6/traccion, 6/7/traccion,
        8/9/compresion, 9/10/compresion, 10/11/compresion, 11/12/compresion,
        1/8/compresion, 12/7/compresion,
        2/8/traccion, 3/9/compresion, 4/10/neutro, 5/11/compresion,
        6/12/traccion,
        8/3/traccion, 9/4/traccion, 11/4/traccion, 12/5/traccion}
      \draw[\tono, line width=4.5pt] (p\i) -- (p\j);
    \foreach \i in {1,...,12}
      \fill[white] (p\i) circle (3.4pt);
    \foreach \i in {1,...,12}
      \draw[principal, line width=1.1pt] (p\i) circle (3.4pt);
  \end{scope}

  % Autor
  \node[anchor=south west, font=\sffamily\Large, text=principal]
    at ([xshift=0.11\paperwidth, yshift=0.08\paperheight]current page.south west)
    {Eddy Eder Sucapuca Cruz};
\end{tikzpicture}
\end{titlepage}

% ---------------------------------------------------------------------------
% Página legal (reverso de la portada)
\thispagestyle{empty}
\null\vfill
{\small\raggedright
\textbf{Análisis matricial de armaduras con Python}\\
Aprende el método de rigidez construyendo tu propia calculadora\par
\medskip
Eddy Eder Sucapuca Cruz\par
Versión preliminar, \the\year.\par
\bigskip
© \the\year{} Eddy Eder Sucapuca Cruz.\par
\medskip
El texto y las figuras de este libro se distribuyen bajo la licencia Creative
Commons Atribución-NoComercial-CompartirIgual 4.0 Internacional (CC BY-NC-SA
4.0): \url{https://creativecommons.org/licenses/by-nc-sa/4.0/deed.es}. Se
permite copiarlos, distribuirlos y adaptarlos sin fines comerciales, citando
al autor y compartiendo las obras derivadas con la misma licencia.\par
\medskip
El código de la calculadora y de los notebooks se distribuye bajo la licencia
MIT: puede usarse, modificarse e incorporarse en otros programas, también
comerciales, conservando el aviso de autoría.\par
\medskip
El código, los notebooks y la versión web del libro están en
\url{https://\textorepositorio}.\par
\bigskip
\textbf{Cómo citar este libro:}\\
Sucapuca Cruz, E. E. (\the\year). \emph{Análisis matricial de armaduras con
Python: aprende el método de rigidez construyendo tu propia calculadora}.
\url{https://\textorepositorio}\par
\bigskip
Compuesto con \LaTeX{} en las tipografías Libertinus y Source Code Pro. Todos
los resultados numéricos y los gráficos del libro se generan ejecutando el
código que se muestra.\par}
\cleardoublepage

\tableofcontents
% capitulos/prefacio.tex
\addchap{Prefacio}

\pendiente{cómo usar este libro (docs/01-indice.md)}


\mainmatter
\part{Primeros pasos}
% capitulos/cap01_por_que_programar.tex
\chapter{Por qué programar el método de rigidez}\label{cap:por_que_programar}

Empecemos por el final. El gráfico siguiente muestra el análisis completo de
un puente peatonal de 18 metros: arriba, cómo se deforma bajo su carga;
abajo, la fuerza axial de cada una de sus 21 barras, en azul las que trabajan
en tracción y en naranja las que trabajan en compresión.

\grafico{c12_analisis_c}

No lo dibujó un programa comercial. Lo dibujó una calculadora de unas
doscientas líneas de Python, escritas una por una a lo largo de este libro.
Al terminarlo, esa calculadora será la nuestra: la habremos escrito, entendido
y verificado paso a paso, y podremos usarla para analizar cualquier armadura
plana.

\begin{objetivos}
  \item Reconocer los límites del cálculo a mano de armaduras.
  \item Explicar qué hace por dentro un programa de análisis estructural.
  \item Describir los pasos del método de rigidez y su orden.
  \item Conocer los tres ejemplos que recorren el libro.
\end{objetivos}

% ---------------------------------------------------------------------------
\section{Los límites del cálculo a mano}\label{sec:limites}\index{método de los nudos}\index{isostática}

En los cursos de estática aprendimos a analizar armaduras con el método de
los nudos y el método de las secciones. Son métodos elegantes y conviene
dominarlos: con ellos verificaremos nuestra calculadora a lo largo del libro.
Pero tienen dos límites.

El primero es el \emph{tamaño}. El puente del gráfico tiene 12 nodos; el
equilibrio de cada nodo da dos ecuaciones, de modo que hay 24 ecuaciones con
24 incógnitas: 21 fuerzas en las barras y 3 reacciones. Resolverlas a mano es
posible, pero largo y propenso a errores, y un error en un nudo se arrastra a
todos los siguientes. Una armadura de cubierta de una nave industrial puede
tener cientos de barras.

El segundo límite es más profundo. Una armadura de $b$ barras, $r$
reacciones y $n$ nodos tiene $b + r$ fuerzas desconocidas y solo $2n$
ecuaciones de equilibrio. Si $b + r > 2n$, la armadura es
\termino{hiperestática}{statically indeterminate}: el equilibrio por sí solo
no basta para obtener sus fuerzas, porque estas dependen además de la
rigidez de las barras. Es el caso de un cuadro con sus dos diagonales, que
estudiaremos como ejemplo B: 6 barras y 3 reacciones frente a 8 ecuaciones.
Ningún método basado solo en el equilibrio puede resolverlo.

El \termino{método de rigidez}{stiffness method} supera ambos límites. Toma
como incógnitas los desplazamientos de los nodos en lugar de las fuerzas,
relaciona cada desplazamiento con las fuerzas que produce mediante la rigidez
de las barras, y reúne todo en un sistema de ecuaciones lineales,
\begin{ecuacionclave}
  \mK\,\vU = \vF,
\end{ecuacionclave}
donde $\mK$ es la \termino{matriz de rigidez}{stiffness matrix} de la
armadura, $\vU$ el vector de desplazamientos de los nodos y $\vF$ el de las
fuerzas aplicadas sobre ellos. El método es el mismo para una armadura
isostática o hiperestática, de tres barras o de trescientas. Y, sobre todo,
es sistemático: cada paso se repite igual para cada barra y cada nodo, lo
que lo hace ideal para un computador.

% ---------------------------------------------------------------------------
\section{Qué hay dentro de un programa comercial}\label{sec:programas}

Los programas de análisis estructural que se usan en la práctica, como
SAP2000, ETABS o Ftool, resuelven las estructuras con el método de rigidez,
generalizado a barras que además se flexionan. Cuando dibujamos una armadura
en uno de ellos, el programa hace lo mismo que haremos en este libro: calcula
la rigidez de cada barra, ensambla la matriz de la estructura, resuelve el
sistema y obtiene las fuerzas de cada barra.

Usar un programa sin entender lo que hace por dentro es riesgoso. El programa
siempre da un resultado, también cuando los datos están mal: un apoyo
olvidado, una carga en la unidad equivocada o un nodo desplazado. Reconocer
que un resultado no tiene sentido exige saber de dónde sale. Por eso este
libro no enseña a usar un programa, sino a construirlo: al escribir cada paso
del método, sabremos qué esperar de él y cómo verificarlo.

\begin{parasabermas}
El método de rigidez se desarrolló en la industria aeronáutica en la década
de 1950, junto con los primeros computadores, y está en la base del método
de los elementos finitos. Su desarrollo completo para armaduras, vigas y
pórticos está en \textcite{kassimali2012} y en \textcite{mcguire2000}.
\end{parasabermas}

% ---------------------------------------------------------------------------
\section{El mapa del método}\label{sec:mapa}\index{método de rigidez!pasos del}

La \cref{fig:mapa_metodo} resume el método de rigidez para armaduras planas
y es, a la vez, el mapa de este libro: cada paso es una parte de nuestra
calculadora, y cada capítulo construye uno o dos pasos.

\begin{figure}[ht]
  \centering
  \figura{mapa_metodo}
  \caption{Los pasos del método de rigidez, con la ecuación principal de
    cada uno y el capítulo donde lo programamos.}
  \label{fig:mapa_metodo}
\end{figure}

En palabras:
\begin{enumerate}
  \item Describimos la armadura con datos: las coordenadas de sus nodos, las
    barras que los unen, los apoyos y las cargas.
  \item Calculamos la longitud y la orientación de cada barra.
  \item Con ellas obtenemos la matriz de rigidez de cada barra, que relaciona
    los desplazamientos de sus extremos con las fuerzas en ellos.
  \item Numeramos los desplazamientos posibles de los nodos, sus
    \termino{grados de libertad}{degrees of freedom}.
  \item Ensamblamos las matrices de las barras en la matriz de la armadura, y
    las cargas en el vector de fuerzas.
  \item Separamos los grados de libertad libres de los restringidos por los
    apoyos.
  \item Resolvemos el sistema: obtenemos los desplazamientos de los nodos y,
    con ellos, las reacciones.
  \item De los desplazamientos obtenemos la fuerza axial de cada barra.
  \item Presentamos, verificamos y dibujamos los resultados.
\end{enumerate}

\begin{recuerda}
Todo el método se basa en tres principios que ya conocemos: el
\emph{equilibrio} de cada nodo, la \emph{compatibilidad} (las barras unidas
en un nodo se desplazan con él) y la \emph{ley de Hooke} (cada barra se
alarga en proporción a su fuerza). Supondremos, como en resistencia de
materiales, que el material es elástico lineal y que los desplazamientos son
pequeños.
\end{recuerda}

% ---------------------------------------------------------------------------
\section{Tres armaduras para aprender}\label{sec:tres_ejemplos}\index{ejemplo A}\index{ejemplo B}\index{ejemplo C}

A lo largo del libro usaremos tres armaduras de acero, con datos realistas
($E = \SI{200}{\giga\pascal}$), que se resumen en la
\cref{tab:tres_ejemplos}. La columna $g$ es el grado de indeterminación,
$g = b + r - 2n$: es cero en una armadura isostática y positivo en una
hiperestática.

\begin{table}[ht]
  \centering
  \caption{Los ejemplos del libro.}
  \label{tab:tres_ejemplos}
  \small
  \begin{tabular}{@{}lcccl@{}}
    \toprule
    Ejemplo & Nodos & Barras & $g$ & Para qué sirve \\
    \midrule
    A. Triángulo & 3 & 3 & 0 & Verificar cada paso a mano \\
    B. Cuadro arriostrado & 4 & 6 & 1 & Un caso hiperestático \\
    C. Puente Pratt & 12 & 21 & 0 & Una armadura de verdad \\
    \bottomrule
  \end{tabular}
\end{table}

El \emph{ejemplo A} es un triángulo de \SI{6}{\metre} de base y
\SI{4}{\metre} de altura, lo bastante pequeño para comprobar a mano cada
número que produzca la calculadora. El \emph{ejemplo B} es el cuadro con dos
diagonales del que hablamos antes: nos recordará que el método resuelve lo
que el equilibrio por sí solo no puede. Y el \emph{ejemplo C} es el puente
del comienzo del capítulo, la prueba final de la calculadora.

Que haya tres ejemplos distintos no es un capricho. Un programa puede dar el
resultado correcto con un ejemplo por casualidad, por ejemplo si confunde dos
números que en ese ejemplo resultan iguales. Probarlo con armaduras de
tamaños y formas distintas hace que esos errores salgan a la luz.

% ---------------------------------------------------------------------------
\section{Lo que necesitamos para empezar}\label{sec:requisitos}

Este libro supone que conocemos la estática de armaduras, las nociones de
esfuerzo y deformación de resistencia de materiales, y el álgebra de
matrices: producto, transpuesta e inversa. No supone ningún conocimiento de
programación. Aprenderemos Python a medida que lo necesitemos, siempre con
ejemplos de armaduras, y cada herramienta nueva aparecerá en una caja como
esta:

\begin{python}
Las cajas de Python explican las herramientas del lenguaje que usamos por
primera vez. El anexo~A las reúne todas para consultarlas rápidamente.
\end{python}

Tampoco necesitamos instalar nada: escribiremos nuestro código en notebooks de
Google Colab, que funcionan en el navegador. En el \cref{cap:primer_notebook}
creamos el primero.

% ---------------------------------------------------------------------------
\repasoyresumen

\begin{itemize}
  \item El cálculo a mano es largo para armaduras grandes e insuficiente para
    las hiperestáticas, cuyo grado de indeterminación $g = b + r - 2n$ es
    positivo.
  \item El método de rigidez toma como incógnitas los desplazamientos de los
    nodos y resuelve el sistema $\mK\,\vU = \vF$. Es el mismo para cualquier
    armadura y es el que usan los programas comerciales.
  \item El método tiene nueve pasos (\cref{fig:mapa_metodo}), que
    programaremos en los capítulos 2 a 12.
  \item Los ejemplos A, B y C, de tamaños y propiedades distintos, nos
    permitirán verificar la calculadora.
\end{itemize}

% ---------------------------------------------------------------------------
\ejercicios

\begin{ejercicio}{1.1}
Calcular el grado de indeterminación de los ejemplos A, B y C, sabiendo que
los tres tienen un apoyo articulado y uno de rodillo. ¿Cuáles pueden
resolverse con el método de los nudos?
\end{ejercicio}

\begin{ejercicio}{1.2}
Cada nodo de una armadura plana puede desplazarse en dos direcciones.
¿Cuántas filas y columnas tiene la matriz $\mK$ de cada ejemplo? ¿Cuántos
números tiene la del ejemplo C?
\end{ejercicio}

\begin{ejercicio}{1.3}
El puente del ejemplo C tiene seis paneles. Para un puente del mismo tipo
con $p$ paneles, deducir cuántos nodos y barras tiene, y comprobar que sigue
siendo isostático. ¿De qué tamaño sería la matriz $\mK$ de un puente de 20
paneles?
\end{ejercicio}

% capitulos/cap02_primer_notebook.tex
\chapter{Nuestro primer notebook}\label{cap:primer_notebook}

\pendiente{ver docs/01-indice.md}


\part{Del modelo a la matriz}
% capitulos/cap03_describir_armadura.tex
\chapter{Describir una armadura}\label{cap:describir_armadura}

Antes de calcular cualquier cosa, nuestra calculadora necesita conocer la
armadura: dónde están sus nodos, qué barras los unen, de qué material y
sección son, cómo está apoyada y qué cargas recibe. En el capítulo anterior
guardamos cada coordenada en una variable suelta; con tres barras eso ya
resulta incómodo, y con cincuenta sería imposible. En este capítulo aprendemos
a organizar toda esa información en estructuras de datos de Python, y
construimos el modelo completo del ejemplo A.

\begin{objetivos}
  \item Identificar los datos que definen una armadura plana.
  \item Usar listas y diccionarios, y elegir entre ellos.
  \item Escribir el modelo de una armadura con diccionarios anidados.
  \item Consultar el modelo y dibujarlo con Matplotlib para verificarlo.
\end{objetivos}

% ---------------------------------------------------------------------------
\section{Los datos de una armadura}\label{sec:datos_armadura}

Recordemos las hipótesis con las que analizamos armaduras: las barras son
rectas y se unen en nodos articulados, las cargas actúan solo en los nodos y
el material es elástico lineal. Con ellas, cada barra trabaja únicamente a
fuerza axial y la armadura queda definida por cuatro grupos de datos:

\begin{description}
  \item[Nodos.] Las coordenadas $(x, y)$ de cada nodo.
  \item[Barras.] El nodo inicial y el nodo final de cada barra, su módulo de
    elasticidad $E$ y el área $A$ de su sección.
  \item[Apoyos.] En qué direcciones está impedido el desplazamiento de cada
    nodo apoyado: un apoyo articulado impide ambas; un rodillo, solo una.
  \item[Cargas.] Las componentes $F_x$ y $F_y$ de las fuerzas aplicadas en los
    nodos.
\end{description}

La \cref{fig:ejemplo_a} y la \cref{tab:ejemplo_a} reúnen estos datos para el
ejemplo A, que nos acompañará en los próximos capítulos.

\begin{figure}[ht]
  \centering
  \figura{ejemplo_a}
  \caption{Ejemplo A: armadura triangular isostática.}
  \label{fig:ejemplo_a}
\end{figure}

\begin{table}[ht]
  \centering
  \caption{Datos del ejemplo A. Todas las barras son de acero
    ($E = \SI{200}{\giga\pascal}$) con un área de \SI{20}{\centi\metre\squared}.}
  \label{tab:ejemplo_a}
  \begin{tabular}{@{}cccc@{}}
    \toprule
    Nodo & Coordenadas (\si{\metre}) & Apoyo & Carga (\si{\kilo\newton}) \\
    \midrule
    1 & $(0, 0)$ & articulado & --- \\
    2 & $(6, 0)$ & rodillo (impide $y$) & --- \\
    3 & $(3, 4)$ & --- & $F_x = 20$, $F_y = -50$ \\
    \bottomrule
  \end{tabular}
  \par\medskip
  \begin{tabular}{@{}ccc@{}}
    \toprule
    Barra & Nodo inicial & Nodo final \\
    \midrule
    1 & 1 & 2 \\
    2 & 2 & 3 \\
    3 & 1 & 3 \\
    \bottomrule
  \end{tabular}
\end{table}

% ---------------------------------------------------------------------------
\section{Listas}\label{sec:listas}

Una \termino{lista}{list} es una colección ordenada de valores, escrita entre
corchetes y separada por comas. Por ejemplo, las coordenadas del nodo~3:
\fragmento{c03_lista}
\salida{c03_lista}

Cada elemento ocupa una posición, llamada \termino{índice}{index}, y se
obtiene escribiendo el índice entre corchetes:
\fragmento{c03_indices}
\salida{c03_indices}

\begin{cuidado}
En Python, las posiciones de una lista se cuentan \emph{desde cero}: el primer
elemento es \py{coordenadas_nodo_3[0]} y el segundo,
\py{coordenadas_nodo_3[1]}. Es una convención de la mayoría de los lenguajes
de programación, pero choca con la numeración de nodos y barras, que en
ingeniería empieza en 1.
\end{cuidado}

Ese desfase se vuelve peligroso cuando usamos una lista para guardar
propiedades de las barras. Si guardamos las áreas de las barras 1, 2 y 3 en
una lista (aquí con un área distinta en la barra 3, para notar la
diferencia), el área de la barra~3 está en la posición~2:
\fragmento{c03_areas}
\salida{c03_areas}

Las listas pueden crecer: el método \py{append} agrega un elemento al final, y
la función \py{len} cuenta los elementos:
\fragmento{c03_append}
\salida{c03_append}

% ---------------------------------------------------------------------------
\section{Diccionarios}\label{sec:diccionarios}

Un \termino{diccionario}{dictionary} guarda pares \emph{clave: valor}, entre
llaves. A diferencia de la lista, cada valor se obtiene por su clave, no por
su posición:
\fragmento{c03_diccionario}
\salida{c03_diccionario}

Las claves pueden ser textos, como \py{"x"} e \py{"y"}, o números. Con números
como claves resolvemos el problema de la sección anterior: el área de la
barra~3 se obtiene con la clave~3, sin traducciones.
\fragmento{c03_claves_numericas}
\salida{c03_claves_numericas}

\begin{recuerda}
Usamos \emph{listas} para colecciones en las que solo importa el orden, y
\emph{diccionarios} cuando cada valor tiene un nombre o un número propio. Los
nodos y las barras tienen número propio: por eso los guardaremos en
diccionarios cuyas claves son esos números, empezando en 1.
\end{recuerda}

% ---------------------------------------------------------------------------
\section{El modelo del ejemplo A}\label{sec:modelo_ejemplo_a}

Un valor de un diccionario puede ser, a su vez, otro diccionario. Con esta
idea, los \termino{diccionarios anidados}{nested dictionaries}, escribimos el
modelo completo. Primero los nodos, cada uno con sus coordenadas:
\fragmento{c03_nodos}

Después las barras, cada una con sus nodos y sus propiedades. Un diccionario
puede ocupar varias líneas: Python sabe que continúa mientras la llave no se
haya cerrado.
\fragmento{c03_barras}

Los apoyos indican, para cada nodo apoyado, si el desplazamiento en $x$ y en
$y$ está impedido. Los valores \py{True} (verdadero) y \py{False} (falso) son
los \termino{booleanos}{booleans} de Python:
\fragmento{c03_apoyos}

Por último, las cargas, con su signo según los ejes globales:
\fragmento{c03_cargas}

\begin{estrategia}
Para escribir el modelo de cualquier armadura:
\begin{enumerate}
  \item Dibujar la armadura, numerar los nodos y las barras desde 1 y elegir
    el sentido de cada barra (nodo inicial y nodo final).
  \item Fijar las unidades (en este libro, \si{\kilo\newton} y \si{\metre}) y
    convertir todos los datos a ellas.
  \item Escribir los cuatro diccionarios: \py{nodos}, \py{barras},
    \py{apoyos} y \py{cargas}.
  \item Verificar el modelo consultándolo y dibujándolo.
\end{enumerate}
\end{estrategia}

Separar los apoyos y las cargas de los nodos tiene una ventaja: los nodos sin
apoyo o sin carga simplemente no aparecen en esos diccionarios, sin entradas
vacías que mantener.

% ---------------------------------------------------------------------------
\section{Consultar el modelo}\label{sec:consultar_modelo}

Para leer un valor anidado encadenamos las claves. Los nodos de la barra~2:
\fragmento{c03_consulta}
\salida{c03_consulta}

Y las coordenadas del nodo inicial de la barra~2, usando el número de nodo
como clave del diccionario \py{nodos}:
\fragmento{c03_consulta_encadenada}
\salida{c03_consulta_encadenada}

Esta consulta encadenada, de la barra a sus nodos y de los nodos a sus
coordenadas, es exactamente la que necesitaremos en el
\cref{cap:geometria} para calcular la geometría de todas las barras.

El operador \py{in} pregunta si una clave existe en un diccionario. El nodo~3
no tiene apoyo, pero sí carga:
\fragmento{c03_in}
\salida{c03_in}

Pedir una clave que no existe, como \py{apoyos[3]}, produce el error
\py{KeyError} y detiene el programa. El método \py{get} devuelve en cambio el
valor \py{None} (\enquote{nada}), o el valor que indiquemos:
\fragmento{c03_get}
\salida{c03_get}

Finalmente, \py{len} también cuenta las entradas de un diccionario:
\fragmento{c03_len}
\salida{c03_len}

% ---------------------------------------------------------------------------
\section{Ver la armadura}\label{sec:ver_armadura}

Un error en las coordenadas o en la conectividad es fácil de cometer y difícil
de detectar leyendo números. La mejor verificación es \emph{ver} la armadura.
Para dibujar usamos \concepto{Matplotlib}, la biblioteca de gráficos más usada
de Python, que importamos con el nombre corto \py{plt}:
\fragmento{c03_import_plt}

\py{plt.subplots()} crea una figura con unos ejes; \py{ax.plot} dibuja una
línea entre los puntos cuyas abscisas y ordenadas recibe en dos listas, y el
texto \py{"k-o"} pide línea negra (\py{k}) continua (\py{-}) con círculos en
los extremos (\py{o}). Dibujamos así cada barra:
\fragmento{c03_dibujo}
\grafico{c03_dibujo}

La instrucción \py{ax.set_aspect("equal")} usa la misma escala en ambos ejes,
para que la armadura no se vea deformada.

\textbf{Verificación.} El dibujo coincide con la \cref{fig:ejemplo_a}: tres
nodos en $(0, 0)$, $(6, 0)$ y $(3, 4)$, unidos por tres barras.

\begin{python}
Escribimos una línea \py{ax.plot} por cada barra. Con tres barras es
aceptable; con cincuenta sería tedioso y propenso a errores. En el
\cref{cap:geometria} aprenderemos a recorrer el diccionario \py{barras} con un
bucle y a dibujar cualquier armadura con unas pocas líneas.
\end{python}

\begin{figure}[ht]
  \centering
  \figura{ejemplo_b}
  \caption{Ejemplo B: cuadro arriostrado hiperestático.}
  \label{fig:ejemplo_b}
\end{figure}

\begin{problemaresuelto}[El modelo del ejemplo B]
Escribir el modelo de datos del ejemplo B (\cref{fig:ejemplo_b}): un cuadro
de \SI{3}{\metre} por \SI{4}{\metre} con sus dos diagonales, articulado en el
nodo~1 y con un rodillo en el nodo~2. Todas las barras son de acero, con un
área de \SI{20}{\centi\metre\squared}. Verificar la cantidad de nodos y de
barras.

\solucion
Seguimos la estrategia de resolución: la numeración ya está en la figura, las
unidades son \si{\kilo\newton} y \si{\metre}, y los diccionarios llevan el
sufijo \py{_b} para no reemplazar los del ejemplo A:
\fragmento{c03_problema_resuelto}
\salida{c03_problema_resuelto}

\textbf{Verificación.} La cantidad de nodos y de barras coincide con la
figura, y la barra~5 es la diagonal que va del nodo~1 al nodo~4.
\end{problemaresuelto}

% ---------------------------------------------------------------------------
\repasoyresumen

\begin{itemize}
  \item Una armadura plana queda definida por sus nodos (coordenadas), sus
    barras (nodos, $E$ y $A$), sus apoyos (direcciones impedidas) y sus cargas
    (componentes en los nodos).
  \item Una lista es una colección ordenada; sus posiciones se cuentan desde
    cero. Un diccionario asocia claves con valores.
  \item El modelo de la calculadora usa cuatro diccionarios: \py{nodos},
    \py{barras}, \py{apoyos} y \py{cargas}, con los números de nodo y de
    barra como claves, empezando en 1.
  \item Los valores anidados se leen encadenando claves entre corchetes,
    de la barra a sus nodos y de los nodos a sus coordenadas.
  \item \py{in} pregunta si una clave existe; \py{get} lee una clave sin
    producir un error si no existe; \py{len} cuenta los elementos.
  \item Todo modelo se verifica dibujándolo con \py{plt.subplots} y
    \py{ax.plot}, con la misma escala en ambos ejes.
\end{itemize}

% ---------------------------------------------------------------------------
\ejercicios

\begin{ejercicio}{3.1}
Usando los diccionarios del ejemplo A, obtener las coordenadas de los nodos
de la barra~3 a partir de \py{barras[3]}, sin escribir los números de nodo a
mano, y calcular la longitud de la barra.
\end{ejercicio}

\begin{ejercicio}{3.2}
Una armadura de cubierta tiene los nodos 1 en $(0, 0)$, 2 en $(4, 0)$, 3 en
$(8, 0)$ y 4 en $(4, 1.5)$, en metros. Sus barras son 1: 1--2, 2: 2--3,
3: 1--4, 4: 4--3 y 5: 2--4, todas de acero; las cuatro primeras tienen un área
de \SI{15}{\centi\metre\squared} y la barra~5, de
\SI{10}{\centi\metre\squared}. Está articulada en el nodo~1, tiene un rodillo
en el nodo~3 y una carga de \SI{30}{\kilo\newton} hacia abajo en el nodo~4.
Escribir su modelo con los nombres \py{nodos_cubierta}, \py{barras_cubierta},
\py{apoyos_cubierta} y \py{cargas_cubierta}.
\end{ejercicio}

\begin{ejercicio}{3.3}
Predecir, sin ejecutarlo, qué muestran \py{apoyos.get(3)} y
\py{3 in cargas, 1 in cargas} para el ejemplo A. Comprobar la predicción en
el notebook.
\end{ejercicio}

\begin{ejercicio}{3.4}
Dibujar la armadura del ejemplo B con Matplotlib y comparar el dibujo con la
\cref{fig:ejemplo_b}.
\end{ejercicio}

\puntodecontrol{03}

% capitulos/cap04_geometria.tex
\chapter{La geometría de todas las barras}\label{cap:geometria}

En el \cref{cap:primer_notebook} calculamos la geometría de una barra, y en el
\cref{cap:describir_armadura} describimos la armadura completa con
diccionarios. Ahora unimos ambas cosas: calculamos la longitud y los cosenos
directores de \emph{todas} las barras, sin importar cuántas sean. La
herramienta que lo hace posible es el bucle, la instrucción que repite un
mismo cálculo para cada elemento de una colección. Con ella, nuestra
calculadora deja de depender del tamaño de la armadura.

\begin{objetivos}
  \item Repetir un cálculo con el bucle \py{for}.
  \item Recorrer un diccionario con \py{.items()} y acumular resultados.
  \item Calcular la geometría de todas las barras y guardarla en un
    diccionario nuevo, sin modificar los datos de entrada.
  \item Presentar los resultados en una tabla y dibujar cualquier armadura.
\end{objetivos}

% ---------------------------------------------------------------------------
\section{El mismo cálculo para cada barra}\label{sec:mismo_calculo}

La geometría de una barra depende solo de las coordenadas de sus dos nodos:
\[
  L = \sqrt{\Delta x^2 + \Delta y^2}, \qquad
  c = \frac{\Delta x}{L}, \qquad s = \frac{\Delta y}{L}.
\]
Para el ejemplo A, con tres barras, podríamos copiar tres veces el cálculo del
\cref{cap:primer_notebook}. Pero una armadura de cubierta típica tiene entre
20 y 50 barras, y la de un puente, varios cientos. Copiar el cálculo sería
tedioso y, sobre todo, una fuente segura de errores: basta con olvidar
cambiar un número en una de las copias.

Lo que necesitamos es escribir el cálculo \emph{una sola vez} y pedirle a
Python que lo repita para cada barra del diccionario \py{barras}.

% ---------------------------------------------------------------------------
\section{El bucle \texttt{for}}\label{sec:for}

Un \termino{bucle}{loop} \py{for} repite un bloque de instrucciones para cada
elemento de una colección. En cada repetición, la variable que sigue a
\py{for} toma el valor del elemento siguiente:
\fragmento{c04_for_lista}
\salida{c04_for_lista}

\begin{python}
El bloque que se repite es el que está \emph{sangrado} (desplazado hacia la
derecha) debajo de la línea del \py{for}, que termina en dos puntos. Python
usa la sangría, normalmente de cuatro espacios, para saber qué instrucciones
pertenecen al bucle. Colab la agrega automáticamente al pulsar Enter después
de los dos puntos; para salir del bloque basta con borrarla.
\end{python}

Si recorremos un diccionario, la variable toma, una por una, sus
\emph{claves}. Con cada clave podemos leer el valor correspondiente:
\fragmento{c04_for_dict}
\salida{c04_for_dict}

Casi siempre necesitamos la clave y el valor a la vez. El método
\py{.items()} entrega ambos en cada repetición, y los recibimos en dos
variables separadas por una coma:
\fragmento{c04_items}
\salida{c04_items}

Un uso muy frecuente de los bucles es \emph{acumular} un resultado. Empezamos
con un valor inicial y, en cada repetición, le sumamos algo. El operador
\py{+=} suma al valor que la variable ya tenía: \py{a += b} equivale a
\py{a = a + b}. Así obtenemos, por ejemplo, la longitud total de las barras:
\fragmento{c04_acumulador}
\salida{c04_acumulador}

% ---------------------------------------------------------------------------
\section{La geometría de todas las barras}\label{sec:geometria_todas}

Ya tenemos todas las piezas. Recorremos las barras; para cada una, buscamos
sus nodos en el diccionario \py{nodos}, calculamos las proyecciones, la
longitud y los cosenos directores, y guardamos el resultado en un diccionario
nuevo, \py{geometria}, con el número de barra como clave:
\fragmento{c04_geometria}

Este código forma parte de nuestra calculadora: es la primera pieza que
funciona para \emph{cualquier} armadura descrita con los diccionarios del
\cref{cap:describir_armadura}.

\begin{cuidado}
Guardamos los resultados en un diccionario nuevo, en lugar de agregarlos a
\py{barras}. Así los datos de entrada nunca cambian: si algo sale mal, podemos
volver a ejecutar el cálculo con la seguridad de partir de los mismos datos.
Mezclar datos y resultados en la misma estructura es una causa frecuente de
errores difíciles de rastrear.
\end{cuidado}

Para presentar los resultados recorremos \py{geometria} e imprimimos una fila
por barra. En las f-strings, el número antes del punto indica el
\emph{ancho} de la columna: \py{:7.3f} reserva siete caracteres con tres
decimales, y \py{:>5} alinea el número de barra a la derecha en cinco
caracteres. Así las columnas quedan alineadas:
\fragmento{c04_tabla}
\salida{c04_tabla}

% ---------------------------------------------------------------------------
\section{Verificación}\label{sec:verificacion_geometria}

Los valores coinciden con los que calculamos a mano en el
\cref{cap:primer_notebook}: la barra~1 es horizontal ($c = 1$, $s = 0$), y las
barras 2 y 3 son las hipotenusas de dos triángulos 3-4-5, con $L = 5$, $s =
0.8$ y $c = \pm 0.6$. El signo de $c$ en la barra~2 es negativo porque va del
nodo~2 al nodo~3, hacia la izquierda.

Una verificación que no depende del ejemplo es que, para toda barra, debe
cumplirse $c^2 + s^2 = 1$:
\fragmento{c04_verificacion}
\salida{c04_verificacion}

Finalmente, comprobamos que los datos de entrada no cambiaron: la barra~1
conserva sus cuatro datos originales y nada más.
\fragmento{c04_datos_intactos}
\salida{c04_datos_intactos}

% ---------------------------------------------------------------------------
\section{Dibujar cualquier armadura}\label{sec:dibujar_cualquier}

En el \cref{cap:describir_armadura} escribimos una línea \py{ax.plot} por
cada barra. Con un bucle, el mismo dibujo funciona para cualquier armadura.
Aprovechamos además para escribir el número de cada barra en su punto medio,
con \py{ax.text(x, y, texto)}:
\fragmento{c04_dibujo}
\grafico{c04_dibujo}

\begin{problemaresuelto}[La geometría del ejemplo B]
Calcular la longitud y los cosenos directores de las seis barras del ejemplo
B (\cref{fig:ejemplo_b}) y verificar los resultados.

\solucion
Aplicamos el mismo bucle a los diccionarios \py{nodos_b} y \py{barras_b} del
\cref{cap:describir_armadura}, y guardamos los resultados en
\py{geometria_b}:
\fragmento{c04_problema_resuelto}
\salida{c04_problema_resuelto}

\textbf{Verificación.} Las barras 1 y 3 son horizontales de \SI{3}{\metre} y
las barras 2 y 4, verticales de \SI{4}{\metre}. Las diagonales 5 y 6 miden
\SI{5}{\metre}; la diagonal~5 sube hacia la derecha ($c > 0$) y la 6, hacia la
izquierda ($c < 0$), como se ve en la figura.
\end{problemaresuelto}

Para resolver el problema copiamos el bucle y cambiamos los nombres de los
diccionarios. Funciona, pero es justamente la clase de repetición que
queríamos evitar. En el \cref{cap:rigidez_barra} aprenderemos a darle un
nombre a un cálculo, con una \emph{función}, para usarlo con cualquier
armadura sin copiarlo.

% ---------------------------------------------------------------------------
\repasoyresumen

\begin{itemize}
  \item El bucle \py{for} repite el bloque sangrado debajo de él para cada
    elemento de una colección.
  \item Al recorrer un diccionario, \py{for} entrega sus claves;
    \py{.items()} entrega la clave y el valor a la vez.
  \item Un acumulador empieza con un valor inicial y crece en cada repetición
    con \py{+=}.
  \item La geometría de todas las barras se guarda en un diccionario nuevo,
    \py{geometria}, con el número de barra como clave. Los datos de entrada no
    se modifican.
  \item En las f-strings, \py{:7.3f} fija el ancho y los decimales de una
    columna, y \py{:>5} alinea a la derecha.
  \item Para toda barra debe cumplirse $c^2 + s^2 = 1$.
\end{itemize}

% ---------------------------------------------------------------------------
\ejercicios

\begin{ejercicio}{4.1}
Modificar el bucle de la geometría para calcular también el ángulo $\theta$
de cada barra del ejemplo A, en grados, y mostrarlo para las tres barras.
\end{ejercicio}

\begin{ejercicio}{4.2}
El peso específico del acero es $\gamma = \SI{78.5}{\kilo\newton\per\metre\cubed}$.
Calcular el peso propio total de las barras del ejemplo A, sabiendo que el
peso de una barra es $\gamma A L$.
\end{ejercicio}

\begin{ejercicio}{4.3}
Calcular la longitud total de las barras del ejemplo B con un acumulador.
\end{ejercicio}

\begin{ejercicio}{4.4}
Dibujar la armadura del ejemplo B con un bucle, numerando sus barras, y
comparar el dibujo con la \cref{fig:ejemplo_b}.
\end{ejercicio}

\puntodecontrol{04}

% capitulos/cap05_rigidez_barra.tex
\chapter{La rigidez de una barra}\label{cap:rigidez_barra}

\pendiente{ver docs/01-indice.md}

% capitulos/cap06_grados_de_libertad.tex
\chapter{Los grados de libertad}\label{cap:grados_de_libertad}

\pendiente{ver docs/01-indice.md}

% capitulos/cap07_ensamblaje.tex
\chapter{Ensamblaje de la matriz global}\label{cap:ensamblaje}

\pendiente{ver docs/01-indice.md}


\part{Resolver y entender}
% capitulos/cap08_particion.tex
\chapter{Cargas y partición del sistema}\label{cap:particion}

\pendiente{ver docs/01-indice.md}

% capitulos/cap09_desplazamientos_reacciones.tex
\chapter{Desplazamientos y reacciones}\label{cap:desplazamientos_reacciones}

\pendiente{ver docs/01-indice.md}

% capitulos/cap10_fuerzas_internas.tex
\chapter{Fuerzas internas}\label{cap:fuerzas_internas}

\pendiente{ver docs/01-indice.md}

% capitulos/cap11_visualizar.tex
\chapter{Visualizar los resultados}\label{cap:visualizar}

\pendiente{ver docs/01-indice.md}

% capitulos/cap12_mi_calculadora.tex
\chapter{Mi calculadora de armaduras}\label{cap:mi_calculadora}

\pendiente{ver docs/01-indice.md}


\appendix
% anexos/anexoA_referencia_python.tex
\chapter{Referencia rápida de Python y NumPy}\label{cap:referencia_python}

\pendiente{ver docs/01-indice.md}

% anexos/anexoB_leer_errores.tex
\chapter{Leer un error de Python}\label{cap:leer_errores}

Cometer errores al programar es inevitable, y no es una mala señal: es parte
del trabajo. Lo que distingue a quien programa con soltura no es que no se
equivoque, sino que sabe leer lo que Python le dice cuando algo falla. Este
anexo enseña a leer esos mensajes con los errores más frecuentes al construir
la calculadora. Todos los mensajes que se muestran son reales: los produjo el
código del libro.

% ---------------------------------------------------------------------------
\section{Cómo se lee un mensaje de error}\label{sec:b_leer}

Supongamos que, al describir una armadura, escribimos por descuido un nodo que
no existe: la barra 2 del ejemplo A va al nodo 4, pero la armadura solo tiene
tres nodos.
\fragmento{b_clave}
Python se detiene y muestra este mensaje:
\salida{b_clave}

El mensaje se llama \termino{traza}{traceback} y se lee \emph{de abajo hacia
arriba}:
\begin{enumerate}
  \item \textbf{La última línea dice qué pasó.} Primero el tipo de error,
    \py{KeyError}, y después el detalle: la clave \py{4} no existe en el
    diccionario.
  \item \textbf{Las líneas anteriores dicen dónde pasó.} Cada bloque que
    empieza con \py{File} es un paso. El primero es nuestra celda; el
    segundo, la línea 10 de la función que llamamos,
    \py{calcular_geometria}, que busca el nodo final de cada barra. Los
    signos \py{~}
    y \py{^} subrayan la parte exacta de la línea que
    falló.
  \item \textbf{El primer bloque es nuestra celda.} Ahí está la línea que
    escribimos nosotros y que desencadenó el error.
\end{enumerate}

Que el error aparezca dentro de una función no significa que la función esté
mal: casi siempre el problema está en los datos que le dimos. Aquí,
\py{calcular_geometria} buscó el nodo final de una barra y no lo encontró,
porque nuestro diccionario \py{barras_mal} usa un nodo inexistente.

\begin{parasabermas}
Colab muestra la misma información con otro formato: agrega colores, marca la
línea que falló con una flecha \py{---->} y nombra las celdas de otra
manera. Se lee igual, de abajo hacia arriba. Colab también puede ofrecer un
botón para que una inteligencia artificial explique el error; es una ayuda
útil, pero conviene aprender a leer el mensaje por nuestra cuenta.
\end{parasabermas}

% ---------------------------------------------------------------------------
\section{Un nombre mal escrito}\label{sec:b_nombre}

El error más frecuente de todos es escribir mal un nombre:
\fragmento{b_nombre}
\salida{b_nombre}
Python no conoce \py{geometira} y, como el nombre se parece a uno que sí
existe, nos sugiere el correcto (\emph{Did you mean}, ¿quisiste decir?).

Un \py{NameError} aparece también cuando usamos una variable o una función
cuya celda no hemos ejecutado: por ejemplo, al abrir de nuevo un notebook,
porque Colab no recuerda nada de la sesión anterior. La solución es ejecutar
las celdas anteriores, o todas con \textsf{Entorno de ejecución > Ejecutar
todas}.

% ---------------------------------------------------------------------------
\section{Errores de escritura}\label{sec:b_escritura}

Algunos errores impiden que Python entienda el código, y aparecen antes de que
se ejecute una sola línea. El más común es olvidar los dos puntos al final de
un \py{for}, un \py{if} o un \py{def}:
\fragmento{b_sintaxis}
\salida{b_sintaxis}
El mensaje no empieza con \py{Traceback}, porque el código no llegó a
ejecutarse: indica la línea y, con el signo \py{^}, el
lugar donde Python esperaba los dos puntos.

El otro error de escritura frecuente es olvidar la sangría del bloque que
sigue a un \py{for}:
\fragmento{b_sangria}
\salida{b_sangria}

% ---------------------------------------------------------------------------
\section{Datos equivocados}\label{sec:b_datos}

Si llamamos a una función con menos argumentos de los que necesita, el
mensaje dice cuál falta. Aquí olvidamos la longitud de la barra:
\fragmento{b_argumentos}
\salida{b_argumentos}

Un \py{IndexError} indica que pedimos una posición que no existe. En la
calculadora casi siempre se debe a confundir la numeración desde 1 con los
índices desde 0. Queremos el desplazamiento vertical del nodo 3, que es el
grado de libertad 6, y lo pedimos con \py{U[6]}:
\fragmento{b_indice}
\salida{b_indice}
\py{U} tiene seis elementos, con índices de 0 a 5. El grado de libertad 6
está en la posición 5, y la forma segura de obtenerla es con la función
\py{indices} del \cref{cap:grados_de_libertad}.

Una división por cero suele indicar dos nodos con las mismas coordenadas: la
barra que los une tiene longitud nula.
\fragmento{b_division}
\salida{b_division}

% ---------------------------------------------------------------------------
\section{Los errores de la calculadora}\label{sec:b_calculadora}

Algunos errores los produce la propia calculadora, a propósito, cuando
detecta datos que no tienen solución. Si a la armadura le falta un apoyo,
$\mK_{ll}$ es singular y \py{resolver_desplazamientos} se detiene con un
mensaje propio (\cref{cap:desplazamientos_reacciones}):
\fragmento{b_inestable}
\salida{b_inestable}
La traza recorre tres pasos: nuestra celda llamó a \py{analizar_armadura},
esta llamó a \py{resolver_desplazamientos}, y ahí se detectó el problema. El
mensaje está en español porque lo escribimos nosotros.

% ---------------------------------------------------------------------------
\section{Los errores que no avisan}\label{sec:b_silenciosos}

Los errores más peligrosos son los que no producen ningún mensaje: una carga
en la unidad equivocada, una coordenada mal escrita o un apoyo en el nodo
equivocado dan resultados, pero resultados falsos. Contra ellos solo sirve la
verificación: comprobar el equilibrio, comparar el dibujo con la armadura y
preguntarse si los números tienen sentido físico, como hicimos a lo largo
del libro. La \cref{tab:errores_frecuentes} del \cref{cap:mi_calculadora}
resume los síntomas más comunes.

\begin{estrategia}
Cuando aparece un error:
\begin{enumerate}
  \item Leer la última línea: el tipo de error y su detalle.
  \item Subir hasta el primer bloque y buscar la línea de nuestra celda que
    lo desencadenó.
  \item Revisar los datos de esa línea: nombres, números de nodo y de barra,
    cantidad de argumentos.
  \item Si el mensaje no basta, imprimir con \py{print} los valores que usa
    la línea, para ver cuál no es el que esperábamos.
  \item Corregir, y ejecutar de nuevo la celda y las que dependen de ella.
\end{enumerate}
\end{estrategia}

% anexos/anexoC_trabajar_en_local.tex
\chapter{Trabajar en local}\label{cap:trabajar_en_local}

\pendiente{ver docs/01-indice.md}

% anexos/anexoD_validacion.tex
% Los valores de SAP2000 vienen de validacion/resultados_sap2000.csv y las
% capturas, de capturas/: ver docs/05-validacion-sap2000.md.
\chapter{Validación con software comercial}\label{cap:validacion}

A lo largo del libro verificamos la calculadora con el cálculo a mano
(ejemplo A), con propiedades que deben cumplirse (simetría, equilibrio) y con
el método de las secciones (ejemplo C). Este anexo agrega una verificación
independiente: comparar sus resultados con los de SAP2000, un programa
comercial de análisis estructural de uso habitual en la práctica, para los
ejemplos B y C.

\avisovalidacion

% ---------------------------------------------------------------------------
\section{Cómo modelar una armadura en un programa comercial}\label{sec:d_modelar}

Los programas comerciales analizan pórticos, cuyas barras resisten flexión.
Para que reproduzcan el modelo de armadura del libro hay que cuidar cinco
cosas:
\begin{enumerate}
  \item \textbf{Barras articuladas.} Liberar los momentos en ambos extremos de
    cada barra, o analizar el modelo como armadura plana, sin giros en los
    nodos. Así cada barra trabaja solo con fuerza axial.
  \item \textbf{Sin peso propio.} Los programas suelen incluir el peso de las
    barras en las cargas; el libro no lo considera, de modo que hay que
    anularlo.
  \item \textbf{Las mismas propiedades.} Un material con
    $E = \SI{200}{\giga\pascal}$ y secciones definidas por su área, sin
    importar su forma.
  \item \textbf{Los mismos apoyos y cargas.} Un apoyo articulado en el nodo 1,
    uno de rodillo con el desplazamiento vertical restringido en el nodo 2
    (ejemplo B) o 7 (ejemplo C), y las cargas aplicadas en los nodos.
  \item \textbf{Unidades coherentes.} Trabajar en kilonewtons y metros, como
    la calculadora, o convertir los resultados antes de compararlos.
\end{enumerate}

SAP2000 dibuja los modelos en el plano $xz$: su eje $z$ es el vertical, que
en el libro llamamos $y$. Sus convenciones de signos coinciden con las del
libro: los desplazamientos y las reacciones son positivos en el sentido de
los ejes, y la fuerza axial $P$ es positiva en tracción.

Los resultados de los dos programas deben coincidir en sus primeras cifras
significativas. Diferencias pequeñas en la última cifra se deben al redondeo;
una diferencia mayor indica un error en uno de los dos modelos, casi siempre
en un apoyo, en las unidades o en el peso propio.

% ---------------------------------------------------------------------------
\section{Ejemplo B}\label{sec:d_ejemplo_b}

El ejemplo B (\cref{fig:ejemplo_b}) es hiperestático, de modo que sus fuerzas
dependen de la rigidez de las barras: es la prueba más exigente para la
calculadora. La \cref{fig:sap_b_modelo} muestra su modelo en SAP2000.

\begin{figure}[ht]
  \centering
  \captura{sap_b_modelo}{modelo del ejemplo B en SAP2000, con los números de
    los nodos y de las barras, los apoyos y las cargas.}
  \caption{El ejemplo B modelado en SAP2000.}
  \label{fig:sap_b_modelo}
\end{figure}

La \cref{tab:d_b_desplazamientos} compara los desplazamientos, y la
\cref{tab:d_b_fuerzas}, las reacciones y las fuerzas axiales.

\begin{table}[ht]
  \centering
  \caption{Ejemplo B: desplazamientos de los nodos, en milímetros.}
  \label{tab:d_b_desplazamientos}
  \small
  \begin{tabular}{@{}ccccc@{}}
    \toprule
    Nodo & $u_x$ calculadora & $u_x$ SAP2000 & $u_y$ calculadora & $u_y$ SAP2000 \\
    \midrule
    1 & \valor{d_b_ux_1} & \valorcomercial{d_b_ux_1} & \valor{d_b_uy_1} & \valorcomercial{d_b_uy_1} \\
    2 & \valor{d_b_ux_2} & \valorcomercial{d_b_ux_2} & \valor{d_b_uy_2} & \valorcomercial{d_b_uy_2} \\
    3 & \valor{d_b_ux_3} & \valorcomercial{d_b_ux_3} & \valor{d_b_uy_3} & \valorcomercial{d_b_uy_3} \\
    4 & \valor{d_b_ux_4} & \valorcomercial{d_b_ux_4} & \valor{d_b_uy_4} & \valorcomercial{d_b_uy_4} \\
    \bottomrule
  \end{tabular}
\end{table}

\begin{table}[ht]
  \centering
  \caption{Ejemplo B: reacciones y fuerzas axiales, en kilonewtons.}
  \label{tab:d_b_fuerzas}
  \small
  \begin{tabular}{@{}lcc@{}}
    \toprule
    Resultado & Calculadora & SAP2000 \\
    \midrule
    Reacción $x$ del nodo 1 & \valor{d_b_r_1} & \valorcomercial{d_b_r_1} \\
    Reacción $y$ del nodo 1 & \valor{d_b_r_2} & \valorcomercial{d_b_r_2} \\
    Reacción $y$ del nodo 2 & \valor{d_b_r_4} & \valorcomercial{d_b_r_4} \\
    \midrule
    Barra 1 & \valor{d_b_n_1} & \valorcomercial{d_b_n_1} \\
    Barra 2 & \valor{d_b_n_2} & \valorcomercial{d_b_n_2} \\
    Barra 3 & \valor{d_b_n_3} & \valorcomercial{d_b_n_3} \\
    Barra 4 & \valor{d_b_n_4} & \valorcomercial{d_b_n_4} \\
    Barra 5 & \valor{d_b_n_5} & \valorcomercial{d_b_n_5} \\
    Barra 6 & \valor{d_b_n_6} & \valorcomercial{d_b_n_6} \\
    \bottomrule
  \end{tabular}
\end{table}

Los resultados coinciden en todas las cifras que muestran las tablas. La
deformada de SAP2000 (\cref{fig:sap_b_deformada}) y su diagrama de fuerzas
axiales (\cref{fig:sap_b_axiales}) son, además, los mismos que dibuja nuestra
calculadora (\cref{fig:resultados_ejemplo_b}).

\begin{figure}[ht]
  \centering
  \captura{sap_b_deformada}{deformada del ejemplo B en SAP2000.}
  \caption{Deformada del ejemplo B en SAP2000.}
  \label{fig:sap_b_deformada}
\end{figure}

\begin{figure}[ht]
  \centering
  \captura{sap_b_axiales}{diagrama de fuerzas axiales del ejemplo B en
    SAP2000, con el valor de cada barra.}
  \caption{Fuerzas axiales del ejemplo B en SAP2000.}
  \label{fig:sap_b_axiales}
\end{figure}

% ---------------------------------------------------------------------------
\section{Ejemplo C}\label{sec:d_ejemplo_c}

El ejemplo C (\cref{fig:ejemplo_c}) es isostático: sus fuerzas dependen solo
del equilibrio, y ya las verificamos con el método de las secciones en el
\cref{cap:mi_calculadora}. Compararlas con SAP2000 confirma que ambos modelos
tienen la misma geometría y las mismas cargas; la flecha, en cambio, depende
de las áreas y del módulo de elasticidad, y verifica las rigideces.

\begin{figure}[ht]
  \centering
  \captura{sap_c_modelo}{modelo del ejemplo C en SAP2000, con los números de
    los nodos y de las barras, los apoyos y las cargas.}
  \caption{El ejemplo C modelado en SAP2000.}
  \label{fig:sap_c_modelo}
\end{figure}

\begin{table}[ht]
  \centering
  \caption{Ejemplo C: flecha y reacciones.}
  \label{tab:d_c_flecha}
  \small
  \begin{tabular}{@{}lcc@{}}
    \toprule
    Resultado & Calculadora & SAP2000 \\
    \midrule
    Desplazamiento $u_y$ del nodo 4 (mm) & \valor{d_c_uy_4} & \valorcomercial{d_c_uy_4} \\
    Reacción $y$ del nodo 1 (kN) & \valor{d_c_r_2} & \valorcomercial{d_c_r_2} \\
    Reacción $y$ del nodo 7 (kN) & \valor{d_c_r_14} & \valorcomercial{d_c_r_14} \\
    \bottomrule
  \end{tabular}
\end{table}

\begin{table}[ht]
  \centering
  \caption{Ejemplo C: fuerzas axiales, en kilonewtons.}
  \label{tab:d_c_fuerzas}
  \small
  \begin{tabular}{@{}cccccc@{}}
    \toprule
    Barra & Calculadora & SAP2000 & Barra & Calculadora & SAP2000 \\
    \midrule
    1 & \valor{d_c_n_1} & \valorcomercial{d_c_n_1} & 12 & \valor{d_c_n_12} & \valorcomercial{d_c_n_12} \\
    2 & \valor{d_c_n_2} & \valorcomercial{d_c_n_2} & 13 & \valor{d_c_n_13} & \valorcomercial{d_c_n_13} \\
    3 & \valor{d_c_n_3} & \valorcomercial{d_c_n_3} & 14 & \valor{d_c_n_14} & \valorcomercial{d_c_n_14} \\
    4 & \valor{d_c_n_4} & \valorcomercial{d_c_n_4} & 15 & \valor{d_c_n_15} & \valorcomercial{d_c_n_15} \\
    5 & \valor{d_c_n_5} & \valorcomercial{d_c_n_5} & 16 & \valor{d_c_n_16} & \valorcomercial{d_c_n_16} \\
    6 & \valor{d_c_n_6} & \valorcomercial{d_c_n_6} & 17 & \valor{d_c_n_17} & \valorcomercial{d_c_n_17} \\
    7 & \valor{d_c_n_7} & \valorcomercial{d_c_n_7} & 18 & \valor{d_c_n_18} & \valorcomercial{d_c_n_18} \\
    8 & \valor{d_c_n_8} & \valorcomercial{d_c_n_8} & 19 & \valor{d_c_n_19} & \valorcomercial{d_c_n_19} \\
    9 & \valor{d_c_n_9} & \valorcomercial{d_c_n_9} & 20 & \valor{d_c_n_20} & \valorcomercial{d_c_n_20} \\
    10 & \valor{d_c_n_10} & \valorcomercial{d_c_n_10} & 21 & \valor{d_c_n_21} & \valorcomercial{d_c_n_21} \\
    11 & \valor{d_c_n_11} & \valorcomercial{d_c_n_11} & & & \\
    \bottomrule
  \end{tabular}
\end{table}

También aquí los resultados coinciden, y el diagrama de fuerzas axiales de
SAP2000 (\cref{fig:sap_c_axiales}) muestra el mismo patrón de tracciones y
compresiones que el de nuestra calculadora en el \cref{cap:mi_calculadora}.

\begin{figure}[ht]
  \centering
  \captura{sap_c_axiales}{diagrama de fuerzas axiales del ejemplo C en
    SAP2000, con el valor de cada barra.}
  \caption{Fuerzas axiales del ejemplo C en SAP2000.}
  \label{fig:sap_c_axiales}
\end{figure}

% ---------------------------------------------------------------------------
\section{Conclusión}\label{sec:d_conclusion}

Para una armadura hiperestática y para una isostática de 21 barras, la
calculadora que construimos en este libro reproduce los desplazamientos, las
reacciones y las fuerzas axiales de un programa comercial. No es casualidad:
ambos resuelven el mismo sistema $\mK\,\vU = \vF$ con el mismo método. La
diferencia es que ahora sabemos qué hace cada paso.

% anexos/anexoE_notacion.tex
\chapter{Notación y convenciones de signos}\label{cap:notacion}

\pendiente{ver docs/01-indice.md}

% anexos/anexoF_soluciones.tex
\chapter{Soluciones de los ejercicios}\label{cap:soluciones}

Cada solución se presenta con el código que la resuelve y su resultado. El
mismo código está en el notebook de soluciones de cada capítulo.

% ---------------------------------------------------------------------------
\section*{Capítulo \ref{cap:primer_notebook}}
\addcontentsline{toc}{section}{Capítulo \ref{cap:primer_notebook}}

\soluciondel{2.1}
\fragmento{s02_ejercicio_1}
\salida{s02_ejercicio_1}
La barra es horizontal y apunta en el sentido positivo de $x$: su longitud es
la diferencia de abscisas, $c = 1$ y $s = 0$.

\soluciondel{2.2}
\fragmento{s02_ejercicio_2}
\salida{s02_ejercicio_2}

\soluciondel{2.3}
Con $E = \SI{200}{\giga\pascal} = \SI{2.0e8}{\kilo\newton\per\metre\squared}$ y
$A = \SI{20}{\centi\metre\squared} = \SI{0.0020}{\metre\squared}$:
\fragmento{s02_ejercicio_3}
\salida{s02_ejercicio_3}

\soluciondel{2.4}
\fragmento{s02_ejercicio_4}
\salida{s02_ejercicio_4}
La barra apunta hacia la izquierda y hacia arriba, de modo que su ángulo está
en el segundo cuadrante: el correcto es el de \py{math.atan2}. El cociente
$\Delta y / \Delta x = 3/(-6)$ es igual al de una barra que apunta hacia la
derecha y hacia abajo, y \py{math.atan} no puede distinguir ambos casos.

\soluciondel{2.5}
\fragmento{s02_ejercicio_5}
\salida{s02_ejercicio_5}
Con \py{math.atan(dy / dx)} la división por $\Delta x = 0$ produce el error
\py{ZeroDivisionError} y el programa se detiene.


\backmatter
\nocite{*}
\printbibliography[title={Bibliografía}]
\printindex

\end{document}
